\documentclass{amsart}
\usepackage{amsmath,amssymb,amsthm,hyperref}
\newtheorem{theorem}{Theorem}
\newtheorem{lemma}{Lemma}
\begin{document}

\section*{Solution}

\subsection*{Statement}
We seek all quadruples of perfect squares $A=x^2$, $B=y^2$, $C=z^2$, $D=t^2$,
with $x,y,z,t$ positive integers (or, more generally, positive rationals, or even
positive reals), such that the four successive products
\[
P_1 = A\cdot A,\quad P_2 = A\cdot A\cdot B,\quad P_3 = A\cdot A\cdot B\cdot C,\quad
P_4 = A\cdot A\cdot B\cdot C\cdot D
\]
satisfy $(P_1,P_2,P_3,P_4)=(A,B,C,D)$. Since $A=x^2$, $B=y^2$, $C=z^2$, $D=t^2$,
in terms of $x,y,z,t$ this is the system
\begin{align}
x^4 &= x^2, \tag{E1}\\
x^4 y^2 &= y^2, \tag{E2}\\
x^4 y^2 z^2 &= z^2, \tag{E3}\\
x^4 y^2 z^2 t^2 &= t^2. \tag{E4}
\end{align}

\subsection*{Result}

\begin{theorem}
Let $x,y,z,t$ be positive real numbers (in particular positive integers or positive
rationals). Then the system \textup{(E1)--(E4)} holds if and only if
\[
x=y=z=1,
\]
while $t$ is \emph{arbitrary}. Equivalently, the complete set of admissible quadruples
of squares is
\[
(A,B,C,D)=(1,\,1,\,1,\,t^2),\qquad t>0 .
\]
In the integer setting the solutions are exactly $(A,B,C,D)=(1,1,1,n^2)$ with
$n\in\mathbb{Z}_{\ge 1}$; in the positive-rational setting they are
$(1,1,1,r^2)$ with $r\in\mathbb{Q}_{>0}$. In particular the only solution in which all
four squares are forced to a single value is the trivial one $(1,1,1,1)$, and every
non-trivial solution is obtained from it by leaving the last square $D=t^2$ free.
\end{theorem}

\begin{proof}
Throughout we use that $x,y,z,t$ are \emph{positive} real numbers. Hence each of
$x^2,y^2,z^2,t^2$ is a positive real number, and in particular each is nonzero, so it
may be cancelled from both sides of an equation; moreover, for a positive real $u$,
the equation $u^2=1$ has the unique solution $u=1$ (the negative root $u=-1$ being
excluded by positivity). We prove the two implications separately.

\medskip
\noindent\textbf{Necessity: every solution satisfies $x=y=z=1$.}

\emph{Step 1: (E1) forces $x=1$.}
Equation (E1) reads $x^4=x^2$, equivalently
\[
x^2\,(x^2-1)=0 .
\]
Since $x>0$, we have $x^2>0$, so the factor $x^2$ is nonzero and we conclude
$x^2-1=0$, i.e.\ $x^2=1$. As $x>0$, this gives $x=1$. Consequently $x^4=1$.

\emph{Step 2: (E2) is automatically consistent.}
Substituting $x=1$ (so $x^4=1$) into (E2) gives $y^2=y^2$, which is an identity. Thus
(E2) imposes no further constraint, but is consistent with $x=1$ for every $y>0$.
(This is expected: after dividing (E2) by the positive quantity $y^2$ one recovers
exactly (E1), so (E2) carries no information beyond (E1).)

\emph{Step 3: (E3) forces $y=1$.}
Substituting $x=1$ into (E3) gives $y^2 z^2=z^2$. Since $z>0$, we have $z^2>0$, so we
may divide by $z^2$ to obtain $y^2=1$. As $y>0$, this gives $y=1$.

\emph{Step 4: (E4) forces $z=1$.}
Substituting $x=1$ and $y=1$ into (E4) gives $z^2 t^2=t^2$. Since $t>0$, we have
$t^2>0$, so we may divide by $t^2$ to obtain $z^2=1$. As $z>0$, this gives $z=1$.

Combining Steps 1--4, any solution of (E1)--(E4) with $x,y,z,t>0$ must satisfy
$x=y=z=1$ (with $t>0$ otherwise unrestricted by these deductions).

\medskip
\noindent\textbf{Sufficiency: $x=y=z=1$ with any $t>0$ is a solution.}
Suppose $x=y=z=1$ and let $t>0$ be arbitrary. Then $x^4=1$ and
\[
P_1=x^4=1=x^2=A,\qquad
P_2=x^4 y^2=1=y^2=B,
\]
\[
P_3=x^4 y^2 z^2=1=z^2=C,\qquad
P_4=x^4 y^2 z^2 t^2=t^2=D .
\]
Thus all four equations (E1)--(E4) hold for every $t>0$. Note that $t$ appears on both
sides of (E4) (it cancels), so no constraint is placed on it once $x=y=z=1$;
the last square $D=t^2$ is therefore free.

\medskip
\noindent\textbf{Conclusion.}
By necessity, every solution has $x=y=z=1$; by sufficiency, each such triple together
with any $t>0$ is indeed a solution. Hence the solution set is exactly
\[
\{(x,y,z,t)\in(0,\infty)^4 : x=y=z=1\},
\]
that is $(A,B,C,D)=(1,1,1,t^2)$ with $t>0$. Restricting $t$ to positive integers
(resp.\ positive rationals) yields the integer (resp.\ rational) solutions stated in
the theorem.
\end{proof}

\subsection*{Remarks}
\begin{itemize}
\item The asymmetry of the answer (the parameters $x,y,z$ are pinned to $1$ while $t$
remains free) reflects the structure of the system: each successive product equation
$P_{k}=$ (the $k$-th square) cancels one more variable, but the \emph{last} square
$D=t^2$ appears on \emph{both} sides of the final equation~(E4) and therefore is not
constrained once the earlier variables equal $1$.
\item In particular there is \emph{no} non-trivial quadruple with all four entries
distinct and $\neq 1$: the only freedom is in $D$.
\item The positivity hypothesis is essential. Over arbitrary reals the equation
$u^2=1$ also admits $u=-1$, and $x^4=x^2$ admits $x=0$ as well; restricting to the
positive squares $A,B,C,D$ (i.e.\ positive $x,y,z,t$) is precisely what makes the
roots unique. Note that even the value of a square such as $A=x^2$ is unaffected by
the sign of $x$, so working with positive $x,y,z,t$ loses no square quadruples.
\end{itemize}

\end{document}
